\documentclass[11pt]{article}
\usepackage[margin=1in]{geometry}
\usepackage{amsmath,amssymb,amsthm,mathtools}
\usepackage{hyperref}

\newtheorem{theorem}{Theorem}
\newtheorem{lemma}[theorem]{Lemma}
\newtheorem{proposition}[theorem]{Proposition}
\theoremstyle{definition}
\newtheorem{definition}[theorem]{Definition}
\newtheorem{example}[theorem]{Example}
\theoremstyle{remark}
\newtheorem{remark}[theorem]{Remark}

\title{Sequences and Limits}
\author{math-foundations / concept 10}
\date{}

\begin{document}
\maketitle

\section{Setting}

Throughout, $(X, d)$ denotes a metric space (concept 09). A
\emph{sequence} in $X$ is a function $x : \mathbb{N} \to X$,
written $(x_n)_{n \in \mathbb{N}}$ or simply $(x_n)$. We
study what it means for such a sequence to settle down toward
a point of $X$, both with and without naming the limit in
advance.

\section{Convergence}

\begin{definition}[convergence]
A sequence $(x_n)$ in $X$ \emph{converges} to $L \in X$ if
\[
\forall \epsilon > 0,\ \exists N \in \mathbb{N},\
\forall n \ge N,\ d(x_n, L) < \epsilon.
\]
We write $x_n \to L$ as $n \to \infty$, or
$\lim_{n \to \infty} x_n = L$.
\end{definition}

\begin{remark}
The integer $N$ is allowed to depend on $\epsilon$. The
content of the definition is: \emph{however small a tolerance
$\epsilon$ you demand, the tail of the sequence eventually
sits inside the open ball of radius $\epsilon$ around $L$,
and stays there.}
\end{remark}

\begin{proposition}[uniqueness of limit]
If $x_n \to L$ and $x_n \to L'$ in $(X, d)$, then $L = L'$.
\end{proposition}

\begin{proof}
Suppose $L \ne L'$ and let $\epsilon = d(L, L')/3 > 0$. Choose
$N_1$ with $d(x_n, L) < \epsilon$ for $n \ge N_1$ and $N_2$
with $d(x_n, L') < \epsilon$ for $n \ge N_2$. For
$n \ge \max(N_1, N_2)$ the triangle inequality gives
$d(L, L') \le d(L, x_n) + d(x_n, L') < 2\epsilon =
\tfrac{2}{3} d(L, L')$, a contradiction.
\end{proof}

\section{The Cauchy property and completeness}

\begin{definition}[Cauchy sequence]
A sequence $(x_n)$ in $X$ is \emph{Cauchy} if
\[
\forall \epsilon > 0,\ \exists N \in \mathbb{N},\
\forall m, n \ge N,\ d(x_m, x_n) < \epsilon.
\]
\end{definition}

The Cauchy condition speaks only of the sequence's internal
behaviour: the terms eventually cluster within any tolerance,
without naming a target.

\begin{definition}[completeness]
A metric space $(X, d)$ is \emph{complete} if every Cauchy
sequence in $X$ converges to a point of $X$.
\end{definition}

\begin{example}
$\mathbb{R}$ with the usual metric is complete; this is one
of the defining features of the real line. By contrast,
$\mathbb{Q}$ is not complete: the sequence of decimal
truncations of $\sqrt{2}$ is Cauchy in $\mathbb{Q}$ but has
no rational limit.
\end{example}

\section{Convergent sequences are Cauchy}

The central theorem of this lesson runs in one direction
unconditionally: convergence always implies the Cauchy
property. The converse requires completeness.

\begin{theorem}[every convergent sequence is Cauchy]
\label{thm:conv-implies-cauchy}
Let $(X, d)$ be a metric space and let $(x_n)$ be a sequence
in $X$. If $(x_n)$ converges, then $(x_n)$ is Cauchy.
\end{theorem}

\begin{proof}
Suppose $x_n \to L$ for some $L \in X$. Fix $\epsilon > 0$.
We must produce $N \in \mathbb{N}$ such that
$d(x_m, x_n) < \epsilon$ whenever $m, n \ge N$.

Apply the definition of convergence with tolerance
$\epsilon/2 > 0$: there exists $N \in \mathbb{N}$ such that
\[
n \ge N \quad \Longrightarrow \quad d(x_n, L) < \frac{\epsilon}{2}.
\]
Now let $m, n \ge N$. By the triangle inequality
(from the metric axioms),
\[
d(x_m, x_n) \le d(x_m, L) + d(L, x_n)
< \frac{\epsilon}{2} + \frac{\epsilon}{2} = \epsilon.
\]
Since $\epsilon > 0$ was arbitrary, the sequence $(x_n)$
satisfies the Cauchy criterion.
\end{proof}

\begin{remark}
The $\epsilon/2$ split is the canonical move: we have two
distances to control ($d(x_m, L)$ and $d(L, x_n)$), each can
be made $< \epsilon/2$, and the triangle inequality combines
them into a bound of $\epsilon$. The same trick appears in
proofs that continuous functions preserve Cauchy sequences,
that uniform limits of continuous functions are continuous,
and that completions exist.
\end{remark}

\section{Bolzano--Weierstrass}

\begin{theorem}[Bolzano--Weierstrass]
Every bounded sequence in $\mathbb{R}$ has a convergent
subsequence.
\end{theorem}

This is the prototypical compactness statement: boundedness
plus closedness in $\mathbb{R}^d$ implies that every sequence
has a convergent subsequence, with the limit lying in the
set. The notion of \emph{sequential compactness} (concept 11
and beyond) is the abstract generalisation.

\section{Why the Cauchy criterion is useful}

In applied work --- and especially in numerical analysis and
machine learning --- we often cannot exhibit a limit in
closed form, but we can monitor distances between successive
iterates. Theorem~\ref{thm:conv-implies-cauchy} together with
completeness of $\mathbb{R}^d$ tells us that
\emph{shrinking pairwise distances suffice}: if the iterates
of an optimisation algorithm satisfy
$\|x_m - x_n\| \to 0$ as $m, n \to \infty$, then they
converge to a point of $\mathbb{R}^d$, regardless of whether
we can name that point. This is the foundation of every
fixed-point convergence proof, including Picard--Lindel\"of
(concept 30) and the construction of strong solutions of
SDEs (concept 33).

\end{document}
